%! Polynomial Functions and Real Zeros — Unit 2, Lesson 2.3 — Homework
\documentclass[10pt]{article}
\usepackage{precalculus-article}
\usepackage{precalculus-boxes}
\newcommand{\TallMath}[1]{$\displaystyle #1\rule[-1.4em]{0pt}{3.2em}$}

\begin{document}

\pageheader{Unit 2, Lesson 2.3}{Homework}

\vspace{0.12in}

\begin{enumerate}[itemsep=10pt]

\item Let $p(x) = (x - 3)(x + 2)(x - 1)$.

\begin{enumerate}[label=(\alph*), itemsep=4pt]
  \item The real zeros of $p$ are \blank{1.2cm} , \blank{1.2cm} ,
        \blank{1.2cm}
  \item The degree of $p$ is \blank{1.2cm}
  \item Written as points, the $x$-intercepts of the graph are
        \blank{1.8cm} , \blank{1.8cm} , \blank{1.8cm}
\end{enumerate}

\item Write each polynomial in factored form.

\begin{enumerate}[label=(\alph*), itemsep=4pt]
  \item Degree 3, zeros $0$, $-4$, and $5$, leading coefficient 1:

        $p(x) = $ \blank{5.5cm}
  \item Degree 3, a zero at $x = 2$ of multiplicity 2 and a zero at
        $x = -1$, leading coefficient 1:

        $p(x) = $ \blank{5.5cm}
\end{enumerate}

\boxguard[18]
\item Let $p(x) = (x + 1)^3(x - 2)^2(x - 5)$. Complete the table, then give
      the degree.

\smallskip
\renewcommand{\arraystretch}{1.5}
\begin{tabular}{|c|c|c|c|}
\hline
\textbf{Zero} & \textbf{Multiplicity} & \textbf{Even or odd?} & \textbf{Crosses or tangent?} \\
\hline
$x = -1$ & \blank{1.2cm} & \blank{1.4cm} & \blank{2.4cm} \\
\hline
$x = 2$  & \blank{1.2cm} & \blank{1.4cm} & \blank{2.4cm} \\
\hline
$x = 5$  & \blank{1.2cm} & \blank{1.4cm} & \blank{2.4cm} \\
\hline
\end{tabular}
\smallskip

Degree of $p$: \blank{1.6cm}

\item Let $p(x) = x^3 - x^2 - 8x + 12$.

\begin{enumerate}[label=(\alph*), itemsep=4pt]
  \item Show that $x = 2$ is a zero of $p$.

        \begin{work}
          p(2) &= (2)^3 - (2)^2 - 8(2) + 12 \\
               &= 8 - 4 - 16 + 12 \\
               &= 0
        \end{work}

        $p(2) = $ \blank{1.6cm} , so $(x - 2)$ \textbf{is} / \textbf{is not}
        a factor of $p(x)$.
  \item In fact $p(x) = (x - 2)^2(x + 3)$. The zero $x = 2$ has multiplicity
        \blank{1.2cm}, so the graph is \blank{2.4cm} to the $x$-axis there,
        while at $x = -3$ the graph \blank{2.4cm} the axis.
\end{enumerate}

\boxguard[16]
\item A polynomial $f$ is given at equally spaced inputs. Fill in the
      difference rows.

\smallskip
\renewcommand{\arraystretch}{1.5}
\begin{tabular}{|l|c|c|c|c|c|c|}
\hline
$x$ & 0 & 1 & 2 & 3 & 4 & 5 \\
\hline
$f(x)$ & 0 & 2 & 12 & 36 & 80 & 150 \\
\hline
1st differences & \multicolumn{6}{c|}{\blank{1.2cm} \ \blank{1.2cm} \ \blank{1.2cm} \ \blank{1.2cm} \ \blank{1.2cm}} \\
\hline
2nd differences & \multicolumn{6}{c|}{\blank{1.2cm} \ \blank{1.2cm} \ \blank{1.2cm} \ \blank{1.2cm}} \\
\hline
3rd differences & \multicolumn{6}{c|}{\blank{1.2cm} \ \blank{1.2cm} \ \blank{1.2cm}} \\
\hline
\end{tabular}
\smallskip

The degree of $f$ is \blank{1.6cm}

\item The graph of a polynomial $p$ is tangent to the $x$-axis at $x = -3$
      and crosses the $x$-axis at $x = 1$. Which of the following could be
      $p(x)$?

\begin{enumerate}[label=\Alph*., itemsep=3pt]
  \item $(x + 3)(x - 1)$
  \item $(x + 3)^2(x - 1)$
  \item $(x + 3)(x - 1)^2$
  \item $(x + 3)^2(x - 1)^2$
\end{enumerate}

\end{enumerate}

\vspace{0.1in}

\boxguard
\begin{extensionbox}
A degree-4 polynomial has only two $x$-intercepts, at $x = -1$ and $x = 3$,
and its graph is tangent to the axis at both. Write a possible formula for
it. Then explain why \textit{no} degree-4 polynomial can be tangent to the
$x$-axis at three different points.

\vspace{0.05in}
\writelines{2}
\end{extensionbox}

\vspace{0.1in}

\boxguard
\begin{notesbox}{Preview of Lesson 2.4}
Every polynomial today had exactly as many real zeros as its degree, once
multiplicities were counted. But $p(x) = x^2 + 1$ has degree 2 and
\textit{no} $x$-intercepts at all. Lesson 2.4,
\textit{Polynomial Functions and Complex Zeros}, explains where the missing
zeros went.
\end{notesbox}

\end{document}
